\chapter{Limitations and concluding remarks on the formulation}
\label{ch:limitations}

\section{Model limitations}
\label{sec:limitations}
The following limitations are inherent to the formulation as posed and accompany any future use of
the model.

\textbf{(1) Physical modelling limits.} The wall is treated as homogeneous, isotropic, linearly
thermoelastic (A3); unfilled PEEK is mildly anisotropic and its glass-transition-modified behaviour
at temperatures approaching or exceeding its operating limits is not captured. Creep and
stress relaxation are excluded, so the formulation describes the short-time-scale steady state and
not hours-long drift. Contact resistance at the interface is neglected (A8). Constant fluid
properties are self-consistent at the baseline but set an upper envelope on thermal load through
bound A-4.

\textbf{(2) Buoyancy.} As discussed in Section~\ref{sec:buoyancy-baseline}, the corrected
length-consistent estimate gives $\mathrm{Ri}=1.7\times10^{-3}$ at the baseline, and the
constant-density closure stands on that estimate. The disposition carries its own guard: the
question re-opens at any parameter-study point with $\mathrm{Ri}>10^{-2}$, the Boussinesq term
remains available in the weak form, and the estimate itself is recomputed from computed fields in
every reported run (regime guard A-1--A-6 audit).

\textbf{(3) Support and geometric ideality.} The pinned supports of A9 are idealizations; any real
compliance at the supports changes both the thermal bowing and the buckling ceiling, and the
clamped comparison is recorded in the audit register. The full-edge pin idealization itself
carries a $0.2\%$ stiffening relative to ideal point pins at the baseline aspect ratio, measured by
the plane-strain Q8 computation of register D-26 ($1.0000$ point pins, $0.998$ full-edge, $-0.045$
clamp, all at the mid-surface). That comparison is a finite-mesh result on a single $24$-element mesh, not a
converged strip solution, so it establishes the pin-versus-clamp distinction only at that level; the
earlier estimate of a few per cent is superseded by that measurement, and no check of this effect
remains outstanding beyond it. Entry-length estimates for the hydrodynamic
and thermal development assume the engineering constants used in
Table~\ref{tab:baseline} and are replaced by computed developing solutions in the future.

\textbf{(4) Reduced wall equation.} Equation \eqref{eq:beam} and the beam--column extension
\eqref{eq:beamcol} are design relations at the reduced (frozen-force) von K\'arm\'an level. The
reduction carries a quantified, one-sided error: at the baseline it overpredicts the
self-consistent deflection by $2.3\%$ ($0.438$ against $0.428\,H_0$), the error stays below
$\approx4\%$ inside the admissible envelope ($N_\theta^c/N_{\mathrm{cr}}\le0.33$), and it grows
rapidly beyond $N_\theta^c/N_{\mathrm{cr}}\approx0.5$, where the escalation criterion of
Chapter~\ref{ch:numerics} requires the consistent force update \eqref{eq:nupdate}; the buckling
margin is unaffected. None of the coupled-model conclusions to be obtained computationally
depend on this accuracy beyond its stated order, and the wall model itself is a verification
target (test V-C3).

\textbf{(5) Literature coverage.} The positioning of Chapter~\ref{ch:contribution} is conditional
on completion of the systematic search (Appendix~B, C-7); several bibliographic records used here
required manual extraction and one, \cite{gupta2024photothermo}, remains bibliographically
incomplete (Appendix~B).

\textbf{(6) Data freeze.} The numerical values of Table~\ref{tab:baseline} that depend on the solid
properties or on $\beta_f$ assume datasheet-representative numbers pending a fixed datasheet
(C-6); only the dimensionless structure, not the point values, is data-independent.

\section{Concluding remarks}
The formulation presented here is complete as a mathematical statement: field equations,
constitutive relations, and boundary data are closed; the plane-strain reduction and the associated
design relations are derived rather than asserted; the dimensionless structure of the regime is
identified from the equations; and a baseline operating point is established arithmetically. The
principal open questions are not mathematical but evidentiary: the bibliography items listed in
Appendix~B, the property freeze, and the construction of the two missing benchmark data sets; the
buoyancy question is closed at the baseline by estimate and guarded by a regime check
(Chapter~\ref{ch:baseline}). The coupled computations themselves are deliberately deferred until these
items close; the isolated solver verifications may proceed under their stated conditions. In that
sense the present document is a research formulation in the strict meaning of the word: everything
stated is either established theory with a source, a derivation reproduced in full, or an explicit
agenda item with named conditions.
